\subsection{مبانی نظری و فرمول‌بندی ریاضی}
فرض کنید $\mathbf{X} \in \mathbb{R}^{B \times L \times D}$ ماتریس حالت‌های پنهان ورودی باشد که در آن $B$ اندازه دسته، $L$ طول دنباله و $D$ بعد پنهان مدل است. پروجکشن‌های اولیه پرس‌وجو، کلید و مقدار با ابعاد سرهای نامتقارن تولید می‌شوند:
\begin{equation}
\mathbf{Q} \in \mathbb{R}^{B \times L \times H_q \times d}, \quad \mathbf{K}_{\text{raw}}, \mathbf{V}_{\text{raw}} \in \mathbb{R}^{B \times N_{\text{kv}} \times H_{\text{kv}} \times d}
\end{equation}
که در آن $H_q$ تعداد سرهای پرس‌وجو، $H_{\text{kv}} \le H_q$ تعداد سرهای فشرده کلید/مقدار و $d = D / H_q$ بعد هر سر توجه است.

به منظور سنتز بازنمایی ویژه هر سر پرس‌وجو، شبکه مسیریاب خطی $\mathbf{W}_{\text{route}} \in \mathbb{R}^{d \times H_{\text{kv}}}$ ماتریس لوجیت‌های مسیر را به ازای هر سر پرس‌وجو تولید می‌کند:
\begin{equation}
\mathbf{R}_{\text{logits}} = \frac{\mathbf{Q} \mathbf{W}_{\text{route}}}{\max(\tau_{\text{router}}, 10^{-4})} \in \mathbb{R}^{B \times L \times H_q \times H_{\text{kv}}}
\end{equation}
سپس توزیع وزن‌های پیوسته مسیریابی با تابع بیشینه هموار نرمال‌سازی می‌شود:
\begin{equation}
\mathbf{W}_{\text{route}} = \operatorname{softmax}(\mathbf{R}_{\text{logits}}, \text{dim}=-1)
\end{equation}
برای هر توکن پرس‌وجو در موقعیت زمانی $t$ و سر پرس‌وجوی $h \in \{1, \dots, H_q\}$، بازنمایی کلید و مقدار سنتز‌شده از ترکیب محدب سرهای فشرده کلید/مقدار به دست می‌آید:
\begin{equation}
\tilde{\mathbf{K}}_{b, t, h} = \sum_{k=1}^{H_{\text{kv}}} W_{\text{route}, b, t, h, k} \cdot \mathbf{K}_{\text{raw}, b, :, k}, \quad \tilde{\mathbf{V}}_{b, t, h} = \sum_{k=1}^{H_{\text{kv}}} W_{\text{route}, b, t, h, k} \cdot \mathbf{V}_{\text{raw}, b, :, k}
\end{equation}
سپس امتیازهای توجه و بافتار نهایی با اعمال ضرب داخلی، ماسک‌های علّی/پدینگ و نرم‌سازی بیشینه هموار استاندارد محاسبه می‌گردد:
\begin{equation}
\mathbf{S}_h = \frac{\mathbf{Q}_h \tilde{\mathbf{K}}_h^\top}{\sqrt{d}}, \quad \mathbf{A}_h = \operatorname{softmax}(\mathbf{S}_h + \mathbf{M}), \quad \mathbf{O}_h = \mathbf{A}_h \tilde{\mathbf{V}}_h
\end{equation}
در نهایت، تانسورهای خروجی کلیه سرها ادغام شده و از پروجکشن خروجی عبور داده می‌شوند:
\begin{equation}
\mathbf{O} = \operatorname{Concat}(\mathbf{O}_1, \dots, \mathbf{O}_{H_q}) \mathbf{W}_O
\end{equation}
